\chapter{Prototype} % Main chapter title

\label{Chapter5} % For referencing the chapter elsewhere, use \ref{Chapter1} 

\lhead{Chapter 5. \emph{Prototype}} % This is for the header on each page - perhaps a shortened title}

\section{Initial Graphs}

%We must settle on a name for these types of graphs
\subsection{Fractional charts}
These types of charts shows the sum of time spent sedentary, standing, and walking. Summing over the three classifications makes it easy to get an overview of the day as a whole. Though the summation gives a great overview, details are lost and it is not possible to pinpoint when each activity occurred during the day. (Name of chart)* can be useful for quickly identifying days with a low amount of activity, before spending time looking at more detailed visualisations. It might also serve as an alert for subjects with a low activity level.

\subsubsection{Pie/box chart}
Our first draft for a pie chart is what we primarily think of when the word pie chart is mentioned. We have discussed if we should display percentage, time or both and how this should be implemented. Putting them next to legends, on the pie chart itself, or having it appear when a section of the pie chart is selected are all possible solutions. The pie chart can be replaced by a box chart. The concept is the same but the circle is replaced by a box, inside are 3 boxes, each representing the different activities.
%Would move this paragraph to another section. Maybe one about experimental charts or something. This does not have that much to do with a pie chart?

Instead of using solid colours the pie chart a set of circles or ``balls'' can be used. Each ball represents a interval of continuous activity. So a small number of large sedentary balls represent few but long periods in a sedentary position. A ball chart visualization has more elements to take in, but help illustrate the subjects activity pattern.

\subsubsection{Symbolic}

Another approach to the same type of graph is using illustrations that intuitively show which activity each illustration represents. The pictures are then scaled so that the size represents the sum of the different types of activity compared to each other.

%Rewrote this section
In our design we use stick figures instead of colours to represent the various activities, a sketch can be seen below (Figure \ref{fig:symbolicPie}). The size of each stick figure would reflect time spent in each activity, i.e. a large walking figure would represent a high amount of hours walking. %SKAL VI HA TELL HE ELLER IKKE, SI NOE OM HVORFOR/HVORFOR IKKE
%We can leave this section here as long as we change the section name (not the subsubsection) to reflect that we are talking about graphs that sum over the different activities. That is what these graphs have in common.

\begin{figure}[h!]
	\centering
		\includegraphics[width=0.5\textwidth]{placeholder.jpg}
		\caption{\footnotesize Rough drafts of a symbolic "pie chart"}
		\label{fig:symbolicPie}
\end{figure}

\subsection{Timeline}
Timeline visualisations are effective at illustrating when various activities occurred during the day. Such a bar can be used to identify points during the day where the subject is in a sedentary position for too long. By looking at multiple days, the user can detect patterns in the day where he needs to be more active.

\subsubsection{Scientific}
A simple yet effective timeline was our first approach. The entire day is represented by a horizontal bar, where each colour corresponds to an activity. A symbol can be used to highlight where sedentary time is higher than a set limit. The size of the bar would have to be quite big in order for a user to inspect some of the shorter time frames, therefore the possibility for zoom functionality should be explored.

Instead of having a continuous scale, a blocked approach can be used. The timeline would be divided into 24 blocks, each block corresponding to an hour. A gradient colour scale would be used to represent the amount of activity inside the hour block. This should make it easy to identify hours in the day where prolonged sedentary positions are present. Giving feedback about specific hours might make it easier to interpret and make use of the chart, because you are alerted to certain hours of the day where you should be more active.

%Perhaps represent the block based one instead?
\begin{figure}[h!]
	\centering
		\includegraphics[width=0.5\textwidth]{placeholder.jpg}
		\caption{\footnotesize Rough drafts of a symbolic "pie chart"}
		\label{fig:simpleTimeline}
\end{figure}

\subsection{Symbolic}
The symbolic approach would be to shape it as a clock. shown by the sketch in figure \ref{fig:clock}. Two clocks would be used to split the 24 with the first clock going from 1 to 12 and the second from 13 to 24. Symbols and distinct background colours can help identify hours of the day easier. 

A darker colour palette could be used during hours of the night. This would deter focus away from these hours and help indicate that a high sedentary time is expected in the time period. 

\begin{figure}[h!]
	\centering
		\includegraphics[width=0.5\textwidth]{placeholder.jpg}
		\caption{\footnotesize Rough drafts of a symbolic "pie chart"}
		\label{fig:clock}
\end{figure}
%Skrevet det litt om for å gjøre det mindre formelt
We came up with a suggestions to represent the passing of time during a day by using an animated timeline and stick figures. The timeline will be drawn in real time while stick figures simultaneously perform the corresponding activities. By displaying the day gradually we hope that the subject will gain a firm understanding of their day. This means that this visualization can not be used to gain a quick overview, but is intended to be used when viewing a day for the first time.

The more motivational approach would be to replace the stick figure with an analogy or metaphor. Instead of a stick figure, a flower could be used. Activity would allow the plant to get sunlight, making it grow. Sedentary positions would make the weather cloudy and the flower would be unaffected.

\subsection{Week overview}
To get an overview of the week as a whole it could be useful to classify each day as either active, not very active or inactive. By calculating the overall activity level and classifying the days into three categories the user can easily see which days he need to be more active and which days the activity level is satisfactory. In our sketch the three different classifications are illustrated by smilies (smiling for active days, and inactive for bad days). 

This chart could also be used ad the top level of an interactive application. Each day could then be clicked to show either a timeline or pie chart. The week overview makes it easy to identify the inactive days, which could then be investigated further by clicking on that day.

\section{Protoype goal/focus}

\section{Prototype implementation}

\section{Findings}
